\subsection{Finitely generated algebras over a field}

In this number, k denotes a field.

Lemma 3. --- Let A be a finitely generated k-algebra and \(\mathfrak{p}_{0} \subset \ldots \subset \mathfrak{p}_{m}\) a maximal chain of prime ideals of A. There exists an integer \(n \geqslant m\), a sequence \((x_{1}, \ldots, x_{n})\) of elements of A, algebraically free over \(k\) (A, IV, p. 4), and such that:

\begin{enumerate}
\def\labelenumi{\alph{enumi})}
\item
  A is integral over the ring \(\mathbf{B} = k[x_1, \ldots, x_n]\) ;
\item
  for every j such that \(0 \leqslant j \leqslant m\), the ideal \(\mathfrak{p}_{j} \cap B\) is generated by the \(x_{k}\) with \(1 \leqslant k \leqslant n - m + j\) .
\end{enumerate}

By the normalization lemma (V, § 3, no. 1, th. 1), there exist an integer \(n \geqslant 0\), a finite sequence \((x_1, ..., x_n)\) of algebraically free elements of A and an increasing sequence \((h(j))_{0 \leqslant j \leqslant m}\) of integers \(\leqslant n\) such that the ideal \(\mathfrak{p}_j \cap \mathbf{B}\) is equal to the prime ideal \(\mathfrak{q}_j\) of B generated by the \(x_k\) with \(1 \leqslant k \leqslant h(j)\) and that A is integral over the ring B. Let \(j\) be an integer such that \(0 \leqslant j < m\). By passing to quotients, one deduces from the canonical injection of B into A an injective homomorphism from \(\mathrm{B} / \mathfrak{q}_j\) to \(\mathrm{A} / \mathfrak{p}_j\) which makes \(\mathrm{A} / \mathfrak{p}_j\) a \((\mathrm{B} / \mathfrak{q}_j)\)-finite algebra. Since the ring \(\mathrm{B} / \mathfrak{q}_j\) is isomorphic to a polynomial algebra in \(n - h(j)\) indeterminates over \(k\), it is integrally closed (V, § 1, no. 3, cor. 3 of prop. 13). By th. 2 of no. 3, we therefore have

\[
1 = \operatorname{ht} \left(\mathfrak {p} _ {j + 1} / \mathfrak {p} _ {j}\right) = \operatorname{ht} \left(\mathfrak {q} _ {j + 1} / \mathfrak {q} _ {j}\right) \geqslant h (j + 1) - h (j).
\]

It follows that we have \(h(j + 1) \leqslant h(j) + 1\) and \(\mathfrak{q}_{j+1} \neq \mathfrak{q}_j\), whence \(h(j + 1) = h(j) + 1\). But we have \(h(m) = n\) since \(\mathfrak{q}_m\) is maximal (V, § 2, no. 1, prop. 1), whence finally \(h(j) = n - m + j\) .

THEOREM 3. — Let A be a finitely generated k-algebra.

\begin{enumerate}
\def\labelenumi{\alph{enumi})}
\item
  For every minimal prime ideal  of A, all maximal chains of prime ideals of A starting at  have length equal to the integer \_.
\item
  The ring A is catenary and its dimension is the supremum of the integers \(\deg. \operatorname{tr}_{k}\kappa(\mathfrak{p})\), where \(\mathfrak{p}\) ranges over the minimal prime ideals of A.
\item
  If A is an integral domain, then all maximal chains of prime ideals of A have the same length, and the dimension of A is the transcendence degree over k of the fraction field of A.
\end{enumerate}

Suppose A is integral and consider a maximal chain \(\mathfrak{p}_0 \subset \ldots \subset \mathfrak{p}_m\) of prime ideals of A. We have \(\mathfrak{p}_0 = 0\) . We then deduce from Lemma 3 the existence of an injective homomorphism \(\varphi : k[\mathrm{X}_1, \ldots, \mathrm{X}_m] \to \mathrm{A}\) of \(k\)-algebras that makes A a finite algebra over \(k[\mathrm{X}_1, \ldots, \mathrm{X}_m]\). Consequently, the transcendence degree over \(k\) of the field of fractions of A is equal to \(m\), whence \(c\) ). The assertion \(a)\) follows from assertion \(c)\) applied to the ring , and assertion \(b)\) is an immediate consequence of \(a)\) and of prop. 5 of § 1, no. 2.

COROLLARY 1. --- Let n be a positive integer. Then

\[
\dim (k [ \mathrm{X} _ {1},..., \mathrm{X} _ {n} ]) = n.
\]

For a \(k\)-algebra A of finite type to have dimension \(n\), it is necessary and sufficient that there exist a \(k\)-homomorphism \(\varphi: k[X_1, ..., X_n] \to A\) making A a finite algebra over \(k[X_1, ..., X_n]\) .

This follows from Th. 3, the normalization lemma (V, § 3, no. 1, Th. 1), and Th. 1, a) of no. 3.

COROLLARY 2. --- Let A be a finitely generated integral domain over k. For every prime ideal  of A, we have

\[
\begin{array}{r l} \mathrm{ht} (\mathfrak {p}) & = \dim (\mathrm{A} _ {\mathfrak {p}}) = \dim (\mathrm{A}) - \dim (\mathrm{A} / \mathfrak {p}) \\ & = \dim (\mathrm{A}) - \deg . \operatorname{tr} _ {k} \kappa (\mathfrak {p}). \end{array}
\]

In particular, we have  for every maximal ideal \_ of A. This follows from th. 3 and remark 4 of § 1, no. 3.

COROLLARY 3. --- Let A be a finitely generated k-algebra and let f be an element of A which belongs to no minimal prime ideal of A (for example an element of A which is not a zero divisor, cf. IV, § 1, no. 1, cor. 3 to prop. 2 and no. 3, cor. 1 to prop. 7). Then one has \(_{f}\).

The map \(\mathfrak{p} \mapsto \mathfrak{p}\mathrm{A}_f\) is a bijection from the set of minimal prime ideals of A onto the set of minimal prime ideals of \(\mathrm{A}_f\) . Moreover, the rings \(\mathrm{A} / \mathfrak{p}\) and \(\mathrm{A}_f / \mathfrak{p}\mathrm{A}_f = (\mathrm{A} / \mathfrak{p})_f\) have the same field of fractions. It therefore suffices to apply Th. 3, b).

COROLLARY 4. --- Let A be a finitely generated k-algebra and  a prime ideal of A. a) For  to be maximal, it is necessary and sufficient that the field of fractions of  be a finite extension of k.

\begin{enumerate}
\def\labelenumi{\alph{enumi})}
\setcounter{enumi}{1}
\tightlist
\item
  Let \(f \in \mathbf{A} - \mathfrak{p}\) ; the ideal \(\mathfrak{p}\) is a maximal ideal of \(\mathbf{A}\) if and only if \(\mathfrak{p}\mathbf{A}_f\) is a maximal ideal of \(\mathbf{A}_f\) .
\end{enumerate}

If  is a maximal ideal of A, then A/p is a field, hence a ring of dimension 0; it is a finitely generated extension of k whose transcendence degree is 0 (th. 3, c)), so it is a finite extension of k. Conversely, if the field of fractions of A/p is a finite extension of k, then \_, so  is maximal. Assertion b) follows from assertion a) in view of the fact that A/p and \_ have the same field of fractions.

The assertion \(a\) ) of cor. 4 is a form of the Nullstellensatz (V, § 3, no. 3, prop. 1).

COROLLARY 5. --- Let A be a finitely generated k-algebra,  a prime ideal of A and \((\mathfrak{p}_{i})_{i\in I}\) the family of minimal prime ideals of A contained in . We have:

\[
\begin{array}{r l} \dim_ {\mathfrak {p}} (\mathrm{A}) & = \sup _ {i \in \mathrm{I}} \dim (\mathrm{A} / \mathfrak {p} _ {i}) \\ & = \dim (\mathrm{A} _ {\mathfrak {p}}) + \dim (\mathrm{A} / \mathfrak {p}) \\ & = \dim (\mathrm{A} _ {\mathfrak {p}}) + \deg . \operatorname{tr} _ {k} \kappa (\mathfrak {p}). \end{array}
\]

We have \(\dim_{\mathfrak{p}}(\mathbf{A}) = \sup_{i\in I}\dim_{\mathfrak{p} / \mathfrak{p}_i}(\mathbf{A} / \mathfrak{p}_i)\) . By § 1, no. 3, prop. 6. But, according to Cor. 3, we have \(\dim_{\mathfrak{p} / \mathfrak{p}_i}(\mathbf{A} / \mathfrak{p}_i) = \dim (\mathbf{A} / \mathfrak{p}_i)\) , whence the first equality. By cor. 2, we have \(\dim (\mathbf{A} / \mathfrak{p}_i) = \dim ((\mathbf{A} / \mathfrak{p}_i)_{\mathfrak{p}}) + \dim (\mathbf{A} / \mathfrak{p})\) . The second equality of the corollary therefore follows from the fact that \(\dim (\mathbf{A}_{\mathfrak{p}}) = \sup_{i\in I}\dim ((\mathbf{A} / \mathfrak{p}_i)_{\mathfrak{p}})\) and the third of Thm. 3.

Thus \(\dim_{\mathfrak{p}}(\mathbf{A})\) is the supremum of the lengths of chains of prime ideals of A contained in \(\mathfrak{p}\).

COROLLARY 6. --- Let A be a finitely generated k-algebra not equal to 0, and let n be an integer . The following conditions are equivalent:

\begin{enumerate}
\def\labelenumi{\alph{enumi})}
\item
  For every \(\mathfrak{p} \in \operatorname{Ass}(\mathbf{A})\) , one has \(\dim(\mathbf{A}/\mathfrak{p}) = n\) .
\item
  Every prime ideal associated with A is minimal and all the irreducible components of Spec(A) are of dimension n.
\item
  There exists an injective \(k\)-homomorphism \(\varphi: k[X_1, ..., X_n] \to A\) making \(A\) a \(k[X_1, ..., X_n]\)-module of finite type and torsion-free.
\end{enumerate}

The equivalence of  and \(a)\) is immediate. Let us show that \(b)\) implies \(a)\). By \(c)\)\(b)\), the ring A has dimension \(n\), and there therefore exists (Corollary 1) a \(k\)-homomorphism \(\varphi : k[X_1, ..., X_n] \to A\) making A a \(k[X_1, ..., X_n]\)-module of finite type. For every prime ideal \(\mathfrak{p} \in \operatorname{Ass}(A)\), the ring \(\mathfrak{p}\) is then integral over \(k[X_1, ..., X_n]\), and one therefore has \(n = \dim(A/\mathfrak{p}) = \dim(k[X_1, ..., X_n]/\varphi^{-1}(\mathfrak{p}))\) by theorem 1, \(a)\), of no. 3, whence \(\varphi^{-1}(\mathfrak{p}) = 0\). The image under the injective homomorphism \(\varphi\) of a nonzero element of \(k[X_1, ..., X_n]\) is not a zero divisor in A (IV, § 1, no. 1, Corollary 3 to Proposition 2), whence \(c)\).

Conversely, suppose condition \(c\) ) is satisfied. For every prime ideal \(\mathfrak{p} \in \operatorname{Ass}(A)\) , the homomorphism \(k[X_1, ..., X_n] \to A/\mathfrak{p}\) deduced from \(\varphi\) is injective (loc. cit.). One therefore has \(\dim(A/\mathfrak{p}) = n\) by cor. 1.

Remark 1. --- By Corollary 5, the conditions , \(a)\), and \(b)\)\(c)\) of Corollary 6 imply that one has \(\dim_{\mathfrak{p}}(\mathbf{A}) = \dim(\mathbf{A})\) for every prime ideal \(\mathfrak{p}\) of \(\mathbf{A}\).

PROPOSITION 4. --- Let A and B be two k-algebras of finite type and \(\rho: \mathbf{A} \to \mathbf{B}\) a homomorphism of algebras. Suppose that A is an integral domain and that the A-module B is torsion-free, and let K denote the field of fractions of A. One has

\[
\dim (\mathbf {B}) = \dim (\mathbf {A}) + \dim (\mathbf {B} \otimes_ {\mathbf {A}} \mathbf {K}).
\]

Suppose first that B is an integral domain. The algebra B \(\otimes_{\mathbf{A}} \mathbf{K}\) is then a ring of fractions of B defined by a multiplicative subset not containing 0. It therefore has as its field of fractions the field of fractions L of B. By th. 3, one has

\[
\dim (\mathbf {B}) = \deg . \operatorname{tr} _ {k} \mathrm{L}, \quad \dim (\mathrm{A}) = \deg . \operatorname{tr} _ {k} \mathrm{K},
\]

\[
\dim (\mathbf {B} \otimes_ {\mathbf {A}} \mathbf {K}) = \deg . \operatorname{tr} _ {\mathbf {K}} \mathbf {L}.
\]

Now one has, by the corollary in A, V, p. 106,

\[
\deg . \operatorname{tr} _ {k} \mathrm{L} = \deg . \operatorname{tr} _ {\mathrm{K}} \mathrm{L} + \deg . \operatorname{tr} _ {k} \mathrm{K},
\]

whence the result in this case.

Let us pass to the general case. Every minimal prime ideal  of B consists of zero divisors in B, hence lies over the ideal 0 of A. It follows that the map \(\mathfrak{p} \mapsto \mathfrak{p}.(\mathbf{B} \otimes_{\mathbf{A}} \mathbf{K})\) is a bijection from the set of minimal prime ideals of B onto the set of minimal prime ideals of \(B \otimes_{A} K\). The proposition therefore follows from the first part of the proof and from prop. 6, c) of § 1, no. 3.

COROLLARY. --- Let \(\rho: A \to B\) be an injective homomorphism of \(k\)-algebras of finite type. One has \(\dim(A) \leqslant \dim(B)\) .

Indeed, let  be a minimal prime ideal of A such that . There exists a prime ideal  of B lying over  (II, § 2, no. 6, Proposition 16). By Proposition 4 applied to \(\dim(A) = \dim(A/p)\) and , one has \(\dim(A) = \dim(A/\mathfrak{p}) \leqslant \dim(B/\mathfrak{q}) \leqslant \dim(B)\), whence the corollary.

Lemma 4. --- Let A and B be two integral domains that are k-algebras, M a torsion-free A-module, and N a torsion-free B-module. If the ring  is an integral domain, then \(_{k}\) is a torsion-free module over \(_{k}\)\(_{k}\).

Let K (resp. L) be the field of fractions of A (resp. B). There exists a set I (resp. J) such that M (resp. N) is isomorphic to a submodule of \(\mathbf{K}^{(\mathbf{I})}\) (resp. \(\mathbf{L}^{(\mathbf{J})}\)). The -module \((\mathbf{A} \otimes_k \mathbf{B})\) is then isomorphic to a submodule of \(M \otimes_k N\), which is isomorphic to \(\mathbf{K}^{(\mathbf{I})} \otimes_k \mathbf{L}^{(\mathbf{J})}\). Since \((\mathbf{K} \otimes_k \mathbf{L})^{(\mathbf{I} \times \mathbf{J})}\) is a ring of fractions of the integral ring \(K \otimes_k L\), it is torsion-free as a module over \(A \otimes_k B\)\(A \otimes_k B\), whence the lemma.

PROPOSITION 5. --- Let \(k'\) be an extension of \(k\), A a \(k\)-algebra of finite type, and B a \(k'\)-algebra of finite type.

\begin{enumerate}
\def\labelenumi{\alph{enumi})}
\tightlist
\item
  The \(k'\) -algebra \(\mathbf{A} \otimes_k \mathbf{B}\) is of finite type and we have
\end{enumerate}

\[
\dim (\mathbf {A} \otimes_ {k} \mathbf {B}) = \dim (\mathbf {A}) + \dim (\mathbf {B}).
\]

\begin{enumerate}
\def\labelenumi{\alph{enumi})}
\setcounter{enumi}{1}
\tightlist
\item
  Let  be a prime ideal of ; let  (resp. ) be the inverse image of \(\otimes_{k}\) in A (resp. B). We have
\end{enumerate}

\[
\dim_ {\mathfrak {r}} (\mathrm{A} \otimes_ {k} \mathrm{B}) = \dim_ {\mathfrak {p}} (\mathrm{A}) + \dim_ {\mathfrak {q}} (\mathrm{B}).
\]

Set \(n = \dim(A)\) and \(m = \dim(B)\) . By cor. 1 to th. 3, there exist injective algebra homomorphisms \(\varphi: k[X_1, ..., X_n] \to A\) and \(\psi: k'[Y_1, ..., Y_m] \to B\)\(A\)\(B\) making respectively A and B into finite algebras over \(k[X_1, ..., X_n]\) and \(k'[Y_1, ..., Y_m]\). The homomorphism \(\varphi \otimes \psi\) is then injective and makes \(A \otimes_k B\) a finite algebra over the \(k'\)-algebra \(k[X_1, ..., X_n] \otimes_k k'[Y_1, ..., Y_m]\), which is identified with \(k'[X_1, ..., X_n, Y_1, ..., Y_m]\). One therefore has \(\dim(A \otimes_k B) = n + m\) by virtue of cor. 1 to th. 3, which proves \(a\) ).

Note that when A and B are integral domains, \(A \otimes_k B\) is a nonzero finitely generated \(k'[\mathrm{X}_{1}, \ldots, \mathrm{X}_{n}, \mathrm{Y}_{1}, \ldots, \mathrm{Y}_{m}]\)-module without torsion by Lemma 4, and hence we have

\[
\dim_ {\mathbf {r}} (\mathrm{A} \otimes_ {k} \mathrm{B}) = n + m = \dim (\mathrm{A}) + \dim (\mathrm{B})
\]

for every prime ideal  of \(\otimes_{k}\) by Remark 1.

Let us now prove \(b\) ). Let \(\mathfrak{r}_{0}\) be a minimal prime ideal of \(\otimes_{k}\) contained in \(\mathfrak{r}\), and denote by \(\mathfrak{p}_{0}\) (resp. \(\mathfrak{q}_{0}\)) the inverse image of \(\mathfrak{r}_{0}\) in A (resp. B). The ring  is isomorphic to a quotient of the ring \((\mathbf{A} \otimes_k \mathbf{B}) / \mathfrak{r}_{0}\)\((\mathbf{A} / \mathfrak{p}_{0}) \otimes_k (\mathbf{B} / \mathfrak{q}_{0})\). We therefore have, according to \(a\),

\[
\dim \left(\left(\mathrm{A} \otimes_ {k} \mathrm{B}\right) / \mathfrak {r} _ {0}\right) \leqslant \dim (\mathrm{A} / \mathfrak {p} _ {0}) + \dim (\mathrm{B} / \mathfrak {q} _ {0}).
\]

Applying cor. 5 of th. 3, we deduce the inequality

\[
\dim_ {\mathfrak {r}} (\mathrm{A} \otimes_ {k} \mathrm{B}) \leqslant \dim_ {\mathfrak {p}} (\mathrm{A}) + \dim_ {\mathfrak {q}} (\mathrm{B}).
\]

Conversely, let  (resp. ) be a minimal prime ideal of A (resp. B) contained in \(\mathfrak{p}_{0}\) (resp. \(q_{0}\)). According to the remark made above, we have

\[
\dim (\mathrm{A} / \mathfrak {p} _ {0}) + \dim (\mathrm{B} / \mathfrak {q} _ {0}) = \dim_ {\overline {{\mathfrak {r}}}} ((\mathrm{A} / \mathfrak {p} _ {0}) \otimes_ {k} (\mathrm{B} / \mathfrak {q} _ {0}))
\]

where \(\overline{\mathbf{r}}\) is the image of  under the canonical surjection \(\mathbf{A} \otimes_{k} \mathbf{B} \to (\mathbf{A}/\mathfrak{p}_{0}) \otimes_{k} (\mathbf{B}/\mathfrak{q}_{0})\). The second member of the preceding equality is evidently less than \(\dim_{\mathfrak{r}}(\mathbf{A} \otimes_{k} \mathbf{B})\). Applying cor. 5 to th. 3, we deduce the inequality

\[
\dim_ {\mathfrak {p}} (\mathrm{A}) + \dim_ {\mathfrak {q}} (\mathrm{B}) \leqslant \dim_ {\mathfrak {r}} (\mathrm{A} \otimes_ {k} \mathrm{B}),
\]

which completes the proof.

COROLLARY. --- Let A be a \(k\)-algebra of finite type, \(k'\) an extension of \(k\), and \(\mathbf{A}'\) the \(k'\)-algebra \(\mathbf{A} \otimes_k k'\).

\begin{enumerate}
\def\labelenumi{\alph{enumi})}
\item
  We have \(\dim (\mathbf{A}^{\prime}) = \dim (\mathbf{A})\).
\item
  Let \(\mathfrak{p}'\) be a prime ideal of \(\mathbf{A}'\) and \(\mathfrak{p}\) its inverse image in \(\mathbf{A}\); one has \(\dim_{\mathfrak{p}'}(\mathbf{A}') = \dim_{\mathfrak{p}}(\mathbf{A})\).
\item
  Let \(\mathfrak{p}'\) be a minimal prime ideal of \(\mathbf{A}'\) and \(\mathfrak{p}\) its inverse image in \(\mathbf{A}\). Then \(\mathfrak{p}\) is minimal and we have \(\dim(\mathbf{A}'/\mathfrak{p}') = \dim(\mathbf{A}/\mathfrak{p})\)
\end{enumerate}

\(a)\) and \(b)\) are deduced from Proposition 5 by taking therein \(\mathbf{B} = k'\). Let us prove \(c)\). The ideal  is minimal (no. 1, Proposition 1) and we have

\[
\dim (\mathrm{A} ^ {\prime} / \mathfrak {p} ^ {\prime}) = \dim_ {\mathfrak {p} ^ {\prime}} (\mathrm{A} ^ {\prime}) = \dim_ {\mathfrak {p}} (\mathrm{A}) = \dim (\mathrm{A} / \mathfrak {p}).
\]

Remark 2. --- Suppose the extension  of  is radicial. Then the canonical map  is a homeomorphism.

Indeed, let \(\mathfrak{p} \in \operatorname{Spec}(\mathbf{A})\). By Lemma 1 of § 1, the space \(f^{-1}(\{\mathfrak{p}\})\) is homeomorphic to \(\operatorname{Spec}(\kappa(\mathfrak{p}) \otimes_k k')\). Now the set \(\alpha\) of nilpotent elements of \(\kappa(\mathfrak{p}) \otimes_k k'\) is a prime ideal (A, V, p. 134, corollary), and the quotient ring \(\kappa(\mathfrak{p}) \otimes_k k')/\mathfrak{a}\), which is an integral domain integral over \(\kappa(\mathfrak{p})\), is a field (A, V, p. 16, cor. 1 and p. 10, Proposition 1). Consequently \(f^{-1}(\{\mathfrak{p}\})\) is reduced to a single element. It follows that the map \(f\) is an increasing bijection from \(\operatorname{Spec}(\mathbf{A}')\) onto \(\operatorname{Spec}(\mathbf{A})\), these two sets being ordered by inclusion of prime ideals, and hence induces a bijection between the irreducible closed subsets of \(\operatorname{Spec}(\mathbf{A})\) and those of \(\operatorname{Spec}(\mathbf{A}')\). Since the closed subsets of \(\operatorname{Spec}(\mathbf{A})\) (resp. \(\operatorname{Spec}(\mathbf{A}')\)) are finite unions of irreducible closed subsets, \(f\) is a homeomorphism.

For a generalization, see Exercise 24, p. 98.

